%% ma: L12-B14-0122
%% lop: 12
%% chuong: 5
%% bai: 14
%% dang: 12.14.32
%% loai: TLN
%% muc: VD
%% dap_an: "2,25"
%% nguon: "Ảnh giáo viên gửi 07/11/2026 (Tư Duy Mở - tuduymo.com, đề tự luyện Phương trình mặt phẳng, đề khoảng cách), Câu 18"
%% kien_thuc: "Bài 14 (hai mặt phẳng song song, khoảng cách từ một điểm đến mặt phẳng)"
%% trang_thai: cho-duyet
%% ly_do: "Dạng 12.14.32 mới đề xuất, chờ giáo viên duyệt."
%% ghi_chu: "Tài liệu gốc không kèm đáp án; đáp án do Claude giải và kiểm bằng sympy."
\begin{ex}
Trong không gian $Oxyz$, cho mặt phẳng $(P)\colon x+2y-2z+1=0$ và mặt phẳng $(Q)\colon ax+by+cz+a-3=0$. Biết rằng khoảng cách giữa hai mặt phẳng bằng $2$. Hãy tính giá trị của biểu thức $T=a^2+b^2+c^2$ (làm tròn kết quả đến hàng phần trăm).
\shortans{2,25}
\loigiai{Khoảng cách giữa hai mặt phẳng xác định nên $(Q)\parallel(P)$: $(a;b;c)=t(1;2;-2)$, $t\ne0$. $(Q)\colon x+2y-2z+\dfrac{t-3}{t}=0$.
$d=\dfrac{\left|1-\frac{t-3}{t}\right|}{3}=\dfrac1{|t|}=2\Rightarrow t^2=\dfrac14$, $T=9t^2=2{,}25$.}
\end{ex}
